\section{Related Work}
\label{sec:related}

\paragraph{Structured pruning of LLMs.}
LLM-Pruner~\citep{llmpruner2023} introduced dependency-aware structural pruning for transformers; Bonsai~\citep{bonsai2025} achieves competitive results using only forward passes; GRASPrune~\citep{graspprune2026} advances the state of the art with global gating that jointly prunes FFN channels and KV head groups.
ShortGPT~\citep{shortgpt2024}, SliceGPT~\citep{slicegpt2024}, and LaCo~\citep{yang2024laco} target depth pruning, while Wanda~\citep{wanda2024} addresses unstructured sparsity; \citet{doingmore2026} revisit the structured-vs-unstructured trade-off under test-time scaling.
\citet{shrestha2026layerpruninglimits} show that layer pruning causes severe degradation on generative reasoning even when classification performance is preserved.
For reasoning models, RESP~\citep{resp2025} uses gradient-based importance on self-generated reasoning traces, RAC~\citep{rac2025} reconstructs chain-of-thought during pruning, SSGR~\citep{ssgr2025} uses self-generated calibration data, and \citet{llms2lrms2026} find width pruning more robust than depth pruning for reasoning but do not study RLVR recovery.
A parallel line recognizes that functional structure should guide pruning: FANG~\citep{fang_2025} groups FFN neurons by activation clusters, FuncNet~\citep{funcnet_2025} decomposes LLMs via ICA, and MPruner~\citep{mpruner2024} uses inter-layer CKA~\citep{kornblith2019cka} for depth pruning in CNNs.
Our work uses \emph{intra-layer} CKA between individual structures at a finer granularity, optimizing for RL recovery capacity rather than immediate accuracy preservation.

\paragraph{Pruning-aware recovery.}
\citet{selfhealing2025} demonstrate that RLVR restores pruned reasoning models to near-original accuracy, significantly outperforming SFT, but test only one pruning method without studying whether the pruning criterion affects recovery.
$P^2$ Law~\citep{p2law2025} establishes scaling laws for post-training loss after pruning.
PASER~\citep{paser2025} optimizes recovery data selection, PAT~\citep{pat2025} jointly considers pruning and fine-tuning objectives, RestoreLCC~\citep{restorelcc_2025} contrastively probes critical attention heads for post-hoc compensation, \citet{boiling_frog_2026} adopt iterative prune-tune loops for gradual compression, and OPSD~\citep{opsd2026} compresses reasoning models \emph{after} RLVR (the reverse of our pipeline).
ARRoL~\citep{xu2026arrol} prunes rollouts during RLVR generation for training efficiency, complementary to our focus on which model structures to prune before training.
These approaches optimize complementary axes (data, training, compensation, schedule); our work isolates which \emph{pruning criterion} enables RL-based recovery.
The Lottery Ticket Hypothesis~\citep{frankle2019lth} asks which subnetworks train effectively from random initialization; our recovery capacity spectrum extends this to \emph{paradigm-conditioned trainability}, where the optimal subnetwork depends on the recovery paradigm (RLVR vs.\ SFT).

\paragraph{RLVR training dynamics.}
\citet{rlsqueezes2026} show that RLVR \emph{compresses} incorrect reasoning trajectories via global exploration, while SFT \emph{expands} correct ones via local imitation.
This asymmetry motivates our finding: gradient importance scoring protects structures that RL would compress, while removing the diverse representations RL exploration requires.
RL training instability has been studied in the RLHF context: \citet{gao2023overoptimization} characterize reward overoptimization, \citet{dong2025rlplus} document repetitive output collapse during extended RL training, and \citet{shah2025comedy} analyze compounding bias from reward estimation.
\citet{mukherjee2025rlsubnetworks} show that RL fine-tuning updates only 5--30\% of parameters, raising the question of whether pruning criteria interact with this active subnetwork.
Our work connects this instability literature to structural pruning, showing that the pruning criterion shapes susceptibility to RL collapse.
